Данные приведены в~таблице.

Короткое слово, уже привязанное неразрывным пробелом, повторно
не~отмечается.

Длинные предлоги оборванными не~выглядят: через некоторое время между
образцами появилось различие около пяти миллиметров, а~после проверки перед
испытанием разница исчезла.

Местоимения предлогами не~являются: их результаты и~его выводы учтены.

Частицы примыкают к~предыдущему слову, ими занимается отдельное правило:
результат отличался бы, автор же указывает.
